\documentclass[12pt]{article}
\usepackage[utf8]{inputenc}
\usepackage[T1]{fontenc}
\usepackage{hyperref}
\title{Tonga Language Resources for K-12 Curriculum Development}
\author{Compiled by Claude Sonnet 5 (claude-sonnet-5) for Andr\'as Kornai}
\date{}
\begin{document}
\maketitle

\section{Introduction}
This dossier covers Tonga (ISO 639-3 toi; Glottocode tong1318), Chitonga, of southern Zambia and northern Zimbabwe, with about 1,300,000 speakers. It must not be confused with tog or toh. The search found grammars but no Bible, modern dictionary or school textbooks. It is a language taught as a subject, but nothing in the search results documents the curriculum materials.

\section{Primary Educational Resources}
None found in search results. Hazel Carter's grammar (below) belongs to a series meant to address the shortage of learning materials for Zambian languages.

\section{Linguistic and Reference Materials}
\begin{itemize}
\item Hazel Carter, \emph{An Outline of Chitonga Grammar} (2002, 96 pp., Bookworld, Zambia; sketch of dialects, orthography, morphology and syntax): \url{https://www.africanbookscollective.com/books/an-outline-of-chitonga-grammar}
\item Publisher page, Gadsden Publishers, Zambia: \url{https://africanbookscollective.com/?p=15987}
\item Cecil Hopgood, \emph{Tonga Grammar} (1940; reprinted 1992 by Zambia Educational Publishing House), bibliographic record: \url{https://glottolog.org/resource/reference/id/28513}
\item Julius Torrend, \emph{A Grammar of the Rhodesian Tonga} (1932), bibliographic record: \url{https://glottolog.org/resource/reference/id/122642}
\item Tonga language (Zambia and Zimbabwe), Wikipedia: \url{https://en.wikipedia.org/wiki/Tonga_language_%28Zambia_and_Zimbabwe%29}
\end{itemize}

\section{Religious and Community Resources}
None found in search results. No Chitonga Bible was identified in this search.

\section{Digital and Audio Resources}
None found in search results.

\section{Curriculum Development Context}
Chitonga is reported to be not fully standardized, with spelling varying between writers. A modern dictionary and graded readers were not found and need a dedicated search of Zambian Ministry of Education and publisher sources. Orthography standardization would be a priority.
\end{document}